\documentclass[12pt]{article}
\usepackage{graphicx} % Required for inserting images
\usepackage[portuguese]{babel}
\usepackage{url}
\usepackage[hidelinks]{hyperref}
\usepackage{csquotes}
\usepackage[backend=biber, style=abnt]{biblatex}
\usepackage{bookmark}
\usepackage{mhchem}
\usepackage{soul}
\addbibresource{Artigo_Valor_de_Uso.bib}

%\usepackage[alf]{abntex2cite}
%\bibliographystyle{abntex2-alf}

\title{Sobre o caráter social dos valores de uso}
\author{Nicholas Funari Voltani}
\date{31 de julho de 2026}

\begin{document}

\maketitle

\begin{abstract}
    Um conhecido resultado d'\textit{O Capital} de Marx é que as mercadorias ``não possuem um só átomo'' de valor, ou seja, que o valor de uma mercadoria pertence à dimensão social, não sendo algo imediatamente empírico. Em contraposição a isso, porém, leituras vulgares costumam atribuir à sua faceta de valor de uso um caráter inteiramente material. Neste trabalho, discute-se como a característica de valor de uso de um objeto, apesar de ser ontologicamente resultado de sua ``corporiedade'' material, trata-se não de uma categoria puramente material, mas também de uma categoria de caráter social.
\end{abstract}


\section{O valor de uso enquanto objeto dado}
A princípio, toda mercadoria é, como Marx o coloca, um mero ``objeto externo, uma coisa que, por meio de suas propriedades, satisfaz necessidades humanas de um tipo qualquer'' \parencite[p. 113]{Marx2017}, aparecendo, à primeira vista, como ``uma coisa óbvia, trivial'' \parencite[p. 146]{Marx2017}. Este caráter fundamental de toda mercadoria — o de satisfação de necessidades humanas — é o que define o que a Economia Política chama de valor de uso.\footnote{``A utilidade de uma coisa faz dela um valor de uso.'' \parencite[p. 114]{Marx2017}.} 

Todo valor de uso satisfaz necessidades humanas através de suas propriedades materiais particulares. Mesmo os assim chamados ``bens imateriais'' não escapam desta descrição geral de mercadorias, e, em particular, de valores de uso. Por exemplo, uma aula online\footnote{Que somente trata-se de uma mercadoria quando é \textit{vendida enquanto mercadoria}, destaque-se. ``A produção capitalista não é apenas produção de mercadoria, mas essencialmente produção de mais-valor. [...] Só é \textit{produtivo} o trabalhador que produz mais-valor para o capitalista ou serve à autovalorização do capital. Se nos for permitido escolher um exemplo fora da esfera da produção material, diremos que um mestre-escola é um trabalhador produtivo se não se limita a trabalhar a cabeça das crianças, mas exige trabalho de si mesmo até o esgotamento, a fim de enriquecer o patrão. Que este último tenha investido seu capital numa fábrica de ensino, em vez de numa fábrica de salsichas, é algo que não altera em nada a questão. Assim, o conceito de \textit{trabalhador produtivo} não implica de modo nenhum apenas uma relação de produção entre atividade e efeito útil, entre trabalhador e produto do trabalho, mas também uma relação de produção especificamente social, surgida historicamente e que cola no trabalhador o rótulo de meio direto de valorização do capital.'' \parencite[p. 578, grifo meu]{Marx2017}.} é, antes de tudo, vibrações no ar advindas da fala de algum professor, captadas e processadas por algum microfone e digitalizada por outros dispositivos eletrônicos; consiste em energia elétrica correndo por cabos de transmissão e \textit{data centers} em algum recôndito da Terra; de ondas eletromagnéticas retransmitidas por antenas e satélites para algum estudante; luzes emitidas por uma tela de computador ou celular etc. Em suma, uma aula online, embora não pareça, pressupõe toda esta materialidade em sua capacidade de satisfação de necessidades humanas. O funcionamento do ``setor de serviços'' requer coisas que já damos por dadas, como algum local físico\footnote{Só porque algumas lojas tenham fechado seus locais físicos não quer que seu estoque deixe de ocupar espaço.}, oxigênio/ar, eletricidade (e, portanto, cabos/fibras óticas) etc., os quais inegavelmente possuem alguma ''materialidade".\footnote{Pode ser uma surpresa a alguns que tudo que está na \textit{cloud}, em verdade, tem seus pés bem fincados na materialidade dos \textit{data centers} em que estão hospedados (e nas quantidades ingentes de eletricidade e água que eles consomem para seu funcionamento), cf. \textcite{Monserrate2022}.} Não é a despeito desta materialidade que estes ``bens e serviços imateriais'' são valores de uso, mas justamente \textit{por conta deles}; são estas particularidades que lhes conferem potencialidades úteis. 

Contudo, a capacidade de satisfação de necessidades humanas, característica dos valores de uso, embora diga respeito à ``corporeidade'' destes objetos, não é algo que se restringe a tais atributos materiais. [com a qual eles venham ao mundo.] Isso até mesmo a teologia cristã apreende: embora Deus proclame que ``toda a árvore, em que há fruto que dê semente, ser-vos-á para mantimento'', ele primeiro ordena Adão e Eva a ``sujeitar'' a terra e ``dominar'' os peixes do mar e as aves do céu.\footnote{Gênesis 1:27-29.} Ou seja, pode até ser verdade que o mundo seja pleno de objetos que \textit{possam} ser de usufruto do ser humano, mas isso não quer dizer que eles assim o sejam \textit{porque} são passíveis de tal usufruto. Em suma, dizer que certos objetos possuem valor de uso \textit{não} quer dizer que sua \textit{raison d'être} natural é a de ser valor de uso. Maçãs não são do jeito que são pura e simplesmente para que seres humanos as comam, e sim o inverso: seres humanos as usufruem (comem, consomem etc.) justamente porque suas propriedades (acidez, doçura, textura etc.) são de seu interesse, satisfazem certas necessidades suas (como fome, desejo etc.). De fato, valores de uso, por sua própria definição, são sempre valores de uso \textit{no que concerne a seres humanos}.\footnote{Talvez seja pertinente destacar que o usufruto de valores de uso não se restringe ao seu consumo \textit{imediato}, mas também de seu uso conjuntamente com outros valores de uso: o óleo de cozinha pode ser um ingrediente importante para a satisfação da fome, mas não é bom para matar a sede.}

Não é à toa que Marx diz que toda coisa útil ``é um conjunto de muitas propriedades e pode, por isso, ser útil sob diversos aspectos'', e que, portanto, ``[d]escobrir esses diversos aspectos e, portanto, as múltiplas formas de uso das coisas é um ato histórico'' \parencite[p. 113]{Marx2017}. Aqui fica claro o papel da ação humana no que tange à ``obtenção'' de valores de uso: não só consiste na descoberta das propriedades das coisas e de como elas podem ser úteis (\textit{ao ser humano}, claro), mas também a de \textit{trazê-las à tona}, ou melhor, de apoderar-se de tais potencialidades de forma que lhes sejam úteis. Ou seja, valores de uso não são coisas que meramente ``estão-aí'' no mundo, à espera de serem colhidas das árvores, e sim coisas \textit{produzidas} pela mão humana, coisas que são produto do \textit{trabalho}.

\section{O pôr teleológico e o trabalho}
Vivendo em um mundo com suas próprias legalidades, às quais submete-se enquanto ser neste/deste mundo, o homem tem de satisfazer suas necessidades enquanto ser natural, o que o força, ao risco de extinção, a apreender estas legalidades e delas valer-se como lhe convier. Sua sobrevivência sempre foi, portanto, mediada pelo ``usufruto'' destas legalidades.\footnote{Isso já é apontado por John Stuart Mill, por exemplo: ``\textit{If we examine any other case of what is called the action of man upon nature, we shall find in like manner that the powers of nature, or in other words the properties of matter, do all the work, when once objects are put into the right position. This one operation, of putting things into fit places for being acted upon by their own internal forces, and by those residing in other natural objects, is all that man does, or can do, with matter. He only moves one thing to or from another. He moves a seed into the ground; and the natural forces of vegetation produce in succession a root, a stem, leaves, flowers, and fruit. [etc.] He has no other means of acting on matter than by moving it. Motion, and resistance tomotion, are the only things which his muscles are constructed for. By muscular contraction he can create a pressure on an outward object, which, if sufficiently powerful, will set it in motion, or if it be already moving, will check or modify or altogether arrest its motion, and he can do no more. But this is enough to have given all the command which mankind have acquired over natural forces immeasurably more powerful than themselves (...). Labour, then, in the physical world, is always and solely employed in putting objects in motion; the properties of matter, the laws of nature, do the rest. The skill and ingenuity of human beings are chiefly exercised in discovering movements, practicable by their powers, and capable of bringing about the effects which they desire. But, while movement is the only effect which man can immediately and directly produce by his muscles, it is not necessary that he should produce directly by them all the movements which he requires. (...)}'' \parencite[p. 48]{Mill1909}. O mérito de Lukács, \textit{vis-à-vis} essa compreensão da Economia Política clássica, é de elucidar que tais legalidades naturais tornam-se legalidades \textit{postas}, i.e. que não é meramente pela \textit{contemplação passiva} que elas devêm meios para fins humanos.} O que o distingue dos demais animais é sua antecipação ideal, e cada vez mais complexa, das legalidades naturais das quais pode valer-se para satisfação de suas necessidades, complexificando não só seu emprego delas, como também da complexificação de suas próprias necessidades.\footnote{``A fim de se apropriar da matéria natural de uma forma útil para sua própria vida, ele [o ser humano] põe em movimento as forças naturais pertencentes à sua corporeidade: seus braços e pernas, cabeça e mãos. Agindo sobre a natureza externa e modificando-a por meio desse movimento, ele modifica, ao mesmo tempo, sua própria natureza. Ele desenvolve as potências que nela jazem latentes e submete o jogo de suas forças a seu próprio domínio.'' \parencite[p. 255]{Marx2017}.} 

Lukács o denomina pôr teleológico\footnote{\textit{Teleologische Setzung}, também traduzido como ``posição teleológica''.}, que consiste na efetivação de uma finalidade, idealizada previamente, cuja efetivação emprega meios mais ou menos adequados a esta finalidade concretamente considerada. Para tal ação, o ser humano converte legalidades naturais em legalidades \textit{postas}: por exemplo, a chuva ainda é, em si, chuva, mas pode agora abastecer tanques, regar plantações etc.\footnote{Vide Lukács: ``Hegel enxerga corretamente o duplo aspecto deste processo, que, por um lado, a posição teleológica [\textit{teleologische Setzung}] 'apenas' faz uso da atividade própria da natureza; por outro lado, que a transformação dessa atividade a faz oposta a si própria. Esta atividade da natureza se torna, portanto, sem transformação ontológico-natural de suas bases, \textit{posta}.'' Dessa forma, o ser humano consegue empregar tais legalidades ``em combinações completamente novas, conferir a elas funções e modos de operar completamente novos''. \parencite[p. 19-20, grifo meu]{Lukacs2018}.} Tal caráter de transformar legalidades naturais em legalidades \textit{postas} consiste meramente na ``mediação de sua subordinação à posição teleológica determinada'' \parencite[p. 20]{Lukacs2018}. 

O resultado do trabalho, um objeto cujo fim último é a satisfação de alguma necessidade humana — do estômago ou da imaginação — , é dito ser um ``valor de uso''. Note-se já que, embora necessariamente consista em algo material, porquanto seja oriundo, em última instância, de legalidades naturais, um valor de uso não é exclusivamente algo natural. Por assim dizer, seu vir-a-ser ideal precede seu vir a ser efetivo: um valor de uso é um produto do trabalho humano justamente porque ele é, desde antes de sua produção, destinado ao consumo humano, ao usufruto que o emprega para sanar alguma necessidade humana; ele somente ``vem ao mundo'' com este propósito.

%%% Caráter histórico

% Ademais, a categoria valor de uso é algo eminentemente \textit{histórico}, não só porque trata-se de características descobertas em certos períodos da história humana, mas porque abarca usos, potências, que perseveram no tempo, ou melhor, cujas potências são conhecidas e preservadas ao longo de gerações humanas. Paralelamente, a complexidade que as necessidades humanas adquirem com o tempo tornam a categoria valor de uso algo mais complexo: ela é histórica e, por consequência, devém algo social para que seja \textit{sabidamente} útil. Um exemplo prático: há uma espécie de peixe no Japão que tem de ser cozinhado minuciosamente para que suas toxinas deixem de ser tóxicas ao (corajoso, ou temerário) consumidor final. A descoberta deste peculiar valor de uso gastronômico certamente é algo histórico --- descoberto em certo período da história humana --- e necessariamente \textit{social}, posto que requer a manipulação deliberada de elementos naturais a fim de que se obtenha uma ``utilidade'' desejada/desejável. 

% Outro exemplo é a enorme variedade de espécies de milho no Peru, resultado de milênios de reprodução cruzada de plantas executada pelos povos e civilizações que ali viveram. Fica claro aqui que não foi o desenrolar da ``história natural'' da evolução dessas espécies vegetais que condicionou o desenvolvimento das civilizações andinas, e sim o contrário: foi o desenvolvimento socio-técnico que induziu tal variabilidade genética em escalas de tempo ``humanas'' (contraposto a escalas de tempo ``naturais'', de milhões de anos)\footnote{Vale bem o que Marx diz: ``Com exceção da indústria extrativa, cujo objeto de trabalho é dado imediatamente pela natureza, tal como a mineração, a caça, a pesca etc. [...], todos os ramos da indústria manipulam um objeto, a matéria-prima, isto é, um objeto de trabalho já filtrado pelo trabalho, ele próprio produto de um trabalho anterior, tal como a semente na agricultura. Animais e plantas, que se costumam considerar como produtos naturais, são, em sua presente forma, não só produtos do trabalho, digamos, do ano anterior, mas o resultado de uma transformação gradual, realizada sob controle humano, ao longo de muitas gerações e mediante o trabalho humano.'' \cite[p. 259]{Marx2017a}. Por outro lado, é evidente que fatores menos facilmente manipuláveis --- como no que tange ao esforço e às escalas de tempo envolvidos, como a geografia local, e a própria evolução natural da fauna e flora locais --- condicionaram mais ``unilateralmente'' o modo de vida desses povos.}. Em geral, pode-se dizer, sim, que a ``evolução natural'' dos hábitats influenciam no caráter social de povos que neles vivem, mas não é sempre o caso; de fato, conforme complexificam-se as relações humanas com a natureza, tende-se a inverter isso, como bem vemos pela atual crise climática.

% Como bem colocado por Duayer,
% \quote{``O pôr a finalidade pressupõe, afirma Lukács, uma apropriação espiritual da realidade orientada pelo fim posto, pois só dessa maneira o resultado do trabalho pode ser algo novo, algo que não emergiria de maneira espontânea dos processos próprios da natureza. No entanto, por contraste, assinala Lukács, o reordenamento dos materiais e processos naturais requerido para que eles possam dar origem ao fim posto exige um conhecimento o mais adequado possível desses objetos e processos, precisamente por convertê-los de legalidades (processos) naturais em legalidades postas. Ao contrário do antropomorfismo próprio da possessão espiritual da realidade condicionada pela finalidade planejada [sic], aqui há de prevalecer o máximo de desantropomorfização, pois a consecução do fim não seria possível sem o conhecimento das propriedades dos objetos e processos envolvidos na transferência das causalidades naturais em causalidades postas.'' \cite[p. 127]{Duayer2023a}}

% Ou seja, é através do ato humano que as ``legalidades naturais'' tornam-se uma ``legalidade'' humanamente útil; em outras palavras, as características naturais dos objetos são condições (ontologicamente) necessárias para que sejam valores de uso, \textit{mas não são condições suficientes}!

% Dessa forma, a categoria valor de uso toma um caráter \textit{histórico}, não só porque trata-se de características descobertas em certos períodos da história humana, mas porque abarca usos e potências que perseveram no tempo, ou melhor, cujo conhecimento destes é conhecido e compartilhado ao longo de gerações humanas. Dessa forma, a complexidade que as necessidades humanas adquirem com o tempo tornam a categoria valor de uso algo mais complexo: ela tornou-se uma categoria eminentemente social (ainda que sempre possua algum suporte material) e, por consequência, devém algo histórico, para que seja \textit{sabidamente} útil ao longo do tempo.\footnote{É nesse sentido que Marx diz que ``Fome é fome, mas a fome que se sacia com carne cozida, comida com garfo e faca, é uma fome diversa da fome que devora carne crua, com mão, unha e dente'' \cite[\textbf{checar página}]{Marx2021}: os valores de uso se complexificam junto das necessidades, e as necessidades se expandem mediante a complexificação dos valores de uso.}


É justamente por tratar-se de algo novo, que não vem-a-ser espontaneamente na natureza, que o valor de uso torna-se algo ``não-somente-natural'' — e isso é não mais que sua alcunha de ``satisfazer necessidades humanas''. Ou seja, ele não só é produzido pelo homem, como deve a ele remeter sua existência enquanto ``valor de uso''; quando não o faz, destrói sua característica de valor de uso — ou seja, torna-se inútil —, por mais que perdure enquanto ser natural.\footnote{Pense-se, por exemplo, que o fim da humanidade significaria, por exemplo, o fim dos valores de uso das casas e dos prédios, por mais que eles permaneçam de pé por séculos após este Armageddon. O que se perderia seria justamente sua capacidade de satisfazer aquelas necessidades humanas --- porque não haveria mais seres humanos com necessidades a se satisfazer!}


\section{A mercadoria e o valor}
A sociedade mercantil aparece, como Marx o coloca, como ``uma enorme coleção de mercadorias'' \parencite[p. 113]{Marx2017}, as quais são, antes de tudo, valores de uso, produtos do trabalho visando a satisfação de necessidades humanas. Tais produtos do trabalho, porém, assumem novas determinações nesta sociedade: são produzidos para satisfazer a necessidade \textit{de outrem}, são produzidos para a \textit{troca}. Como Smith famosamente o coloca, ``não é da benevolência do açougueiro, do cervejeiro ou do padeiro, que esperamos nosso jantar''\footnote{``\textit{It is not from the benevolence of the butcher, the brewer, or the baker, that we expect our dinner, but from their regard to their own interest. We address ourselves, not to their humanity but to their self-love, and never talk to them of our own necessities but of their advantages.}'' \parencite[p. 43]{Smith1904}.}, e sim porque tem quem queira comprar seus produtos, os quais eles desejam vender justamente porque desejam comprar de outrem. 

Embora o caráter de valor de uso seja ontologicamente prioritário a esta — e toda — sociedade, aqui é o \textit{valor} que assume momento predominante sobre sua dinâmica: produz-se para vender, e não só para produzir valor, e sim mais-valor\footnote{``[O capitalista] quer produzir não só um valor de uso, mas uma mercadoria; não só valor de uso, mas valor, e não só valor, mas também mais-valor.'' \parencite[p. 263]{Marx2017}.}. Marx, assim como a Economia Política clássica, desvela que a substância do valor é o trabalho abstrato, o qual é algo ``suprassensível'' (``não há um só átomo'' etc.) pertencente à mercadoria, isto é, é algo de caráter social que lhe pertence.

A mercadoria é unidade de valor de uso e valor: somente afirma-se como um quando nega-se como outro, mas não é mercadoria sendo exclusivamente um destes. Aqui, porém, já pressupõe-se que se trata de \textit{frutos do trabalho}, ou seja, já são fruto de uma malha social suficientemente desenvolvida. O fato de um destes produtos ser considerado ``inútil'' ganha novas determinações: valores de uso perfeitamente passíveis de satisfação de necessidades humanas somente são \textit{úteis}, em uma sociedade capitalista, porquanto sejam passíveis de \textit{venda}. Um exemplo infame desta determinação da sociedade capitalista foi a política de queima de dezenas de milhões de sacas de café durante todo o ``primeiro'' governo Vargas, iniciada em 1931 — dois anos depois do crash da Bolsa de Nova York em 1929 — e encerrada somente em 1944, como política extrema de controle dos preços do café: a dependência extrema da economia brasileira dos produtores de café fazia com que a queda de seus preços tornasse sua produção ``inútil'', trazendo o perigo do colapso da economia inteira.\footnote{Num cenário mais factível, porém, tais destruições não-materiais de certos valores de uso também podem ocorrer, por exemplo, quando surgem substitutos de seu uso: \textcite{Foster2004} dão o exemplo da decadência da economia andina após a obsolescência dos nitratos chilenos — principal motor desta economia ao final do século XIX — por conta do desenvolvimento do chamado ``processo de Haber'' no começo do século XX, através do qual produz-se amônia através da fixação do nitrogênio do próprio ar ($\ce{ N_{2} + 3H_{2} <=> 2 NH_{3}}$). ``\textit{The result within a few years was to destroy almost completely the value of Chilean nitrates, creating a severe crisis for the Chilean economy}.'' \parencite[p. 192]{Foster2004}. Embora não tivesse destruído os nitratos em si, destruiu suas chances de venda, portanto tornando sua produção — \textit{no modo de produção capitalista} — imprestável. }

\section{Considerações finais}
Uma consequência fundamental da contradição entre valor de uso e valor é que toda mercadoria apareça como mero objeto de usufruto, i.e. como mero valor de uso, e que seu valor apareça como algo que lhe é ``impresso'' externamente pelo dinheiro, o representante universal de valor no mundo das mercadorias.

[A desenvolver.]



\printbibliography

\end{document}